% !TEX root = ../notes_template.tex
\chapter{Pelvis \& Hip}\label{chp:pelvis_hip}

\section*{Objectives}

\begin{itemize}
  \item Describe the role of the pelvis in orienting the acetabulum and supporting femoral motion during closed and open chain tasks.
  \item Interpret SFMA Top Tier patterns (MSF, MSE, MSR, SLS, ADS) to identify hip-related movement dysfunction.
  \item Perform MMT for primary hip movements (flexion, extension, abduction, adduction, internal rotation, external rotation) with correct patient positioning and stabilization.
  \item Identify muscle group innervation and root levels using neuromechanical reasoning.
  \item Conduct muscle length testing for the hamstrings, iliopsoas, rectus femoris, and TFL/ITB using standardized protocols.
  \item Acquire and interpret RUSI images of the femoral neurovascular bundle.
  \item Classify observed findings using R1/R2 reasoning, mobility/motor control logic, and functional movement integration.
\end{itemize}

\section*{Functional Integration: Pelvis \& Hip as the Proximal Engine}

The pelvis serves as the dynamic base for femoral motion, functioning analogously to the scapula in the upper extremity. Just as the scapula positions the glenoid for effective humeral mechanics, the pelvis orients the acetabulum and stabilizes the proximal femur for forceful, efficient movement.

Movements of the pelvis are, in fact, movements of the spine—primarily the lumbar spine, though contributions from the thoracic and sacral regions may be relevant depending on posture and load. These pelvic motions (e.g., anterior tilt, rotation) will be addressed in Week 8 during the Lumbopelvic Anatomy lab.

It is important to recognize that hip movements, particularly in closed chain positions, are often coupled with pelvic and spinal motions due to kinetic chain continuity. For example, an anterior pelvic tilt (achieved with lumbar extension without changing the position of the thorax in space) in standing cannot occur without concurrent hip flexion. This coupling reflects the interdependence of segmental mobility and neuromuscular control, and underscores the need for integrated assessment strategies when evaluating hip function in functional tasks.

Hip function cannot be fully understood in isolation—it is deeply integrated with lumbopelvic control and distal alignment. Dysfunction at the hip may produce lumbar pain, pelvic asymmetry, or gait deviation, while impaired trunk or knee mechanics often constrain hip mobility.

\begin{tcolorbox}[colback=gray!5, colframe=black, title={Functional Movement Approach: Top Tier SFMA Patterns}]
SFMA patterns such as MSF, MSE, MSR, and SLS rely on coordinated pelvis–hip action. Deficits in these tests may reflect mobility limitations, length impairments, or poor motor control rooted in the lumbopelvic–hip complex. The Arms Down Squat (ADS) adds an integrated assessment of hip, knee, and ankle synergy.
\end{tcolorbox}

This week’s lab uses these patterns as entry points for deeper evaluation through manual muscle testing, muscle length assessment, and RUSI-guided observation. Our aim is to link functional movement testing to segmental neuromechanical reasoning.

\subsection*{SFMA Top Tier Patterns for Pelvic-Hip Focus}

\begin{itemize}
  \item \textbf{Multisegmental Flexion (MSF)} — Observe symmetry and depth of hip hinge; limited flexion may stem from posterior chain tightness, altered lumbo-pelvic rhythm, or hip capsule restriction.
  \item \textbf{Multisegmental Extension (MSE)} — Assess anterior hip and trunk extension; pelvic thrust or asymmetry may indicate impaired hip extension, rectus femoris tightness, or poor gluteal activation.
  \item \textbf{Multisegmental Rotation (MSR)} — View pelvic dissociation and hip rotation bias; asymmetries here may flag lumbar-pelvic-hip mobility restrictions.
  \item \textbf{Single Leg Stance (SLS)} — Detect frontal plane pelvic control; Trendelenburg signs or instability may point to abductor weakness, proprioceptive deficits, or gluteal inhibition.
  \item \textbf{Arms Down Squat (ADS)} — Assess hip, knee, and ankle coordination with arms relaxed; limited depth or pelvic asymmetry may suggest hip joint restriction, motor control deficits, or compensatory strategies from above or below.
\end{itemize}

\vspace{0.25cm}

\noindent
In practice clinicians often use functional movement observations (SFMA being one example) to guide which structures to assess further via MMT, MLT, and neurovascular screening. Consider whether findings are mobility- or motor-control-driven to apply the appropriate intervention focus.


\section*{Hip Movements: Muscles and Innervation}

\vspace{0.25cm}

\noindent
The table below lists primary hip movements with corresponding prime movers, peripheral nerves, and spinal roots. Movements are organized top-down by functional direction, with muscles listed from proximal (higher root) to distal (lower root) to aid in neurosegmental reasoning.

\begin{table}[H]
\centering
\captionsetup{justification=raggedright, singlelinecheck=false}
\begin{tabular}{|p{3.5cm}|p{5cm}|p{4cm}|p{2.5cm}|}
\hline
\textbf{Movement} & \textbf{Muscle} & \textbf{Peripheral Nerve} & \textbf{Nerve Root(s)} \\
\hline
\multirow{3}{*}{Hip Flexion}
  & Iliopsoas & Femoral / Rami (L1–L3) & L1–L3 \\
  & Rectus Femoris & Femoral & L2–L4 \\
  & Sartorius & Femoral & L2–L3 \\
\hline
\multirow{3}{*}{Hip Extension}
  & Gluteus Maximus & Inferior Gluteal & L5–S2 \\
  & Hamstrings & Sciatic & L5–S2 \\
\hline
\multirow{2}{*}{Hip Abduction}
  & Gluteus Medius & Superior Gluteal & L4–S1 \\
  & Gluteus Minimus & Superior Gluteal & L4–S1 \\
\hline
\multirow{2}{*}{Hip Adduction}
  & Adductor Longus & Obturator & L2–L4 \\
  & Adductor Magnus & Obturator / Sciatic & L2–L4 \\
\hline
\multirow{2}{*}{External Rotation}
  & Piriformis & Nerve to Piriformis & L5–S2 \\
  & Obturator Externus & Obturator & L3–L4 \\
\hline
\multirow{2}{*}{Internal Rotation}
  & Tensor Fascia Latae & Superior Gluteal & L4–S1 \\
  & Gluteus Medius (ant) & Superior Gluteal & L4–S1 \\
\hline
\end{tabular}
\caption{Primary Hip Movements and Their Neuromuscular Innervation}
\end{table}

\subsection*{MMT Hip}

\begin{labactivitybox}
\textbf{Activity: MMT – Hip Flexion and Extension}

\begin{itemize}
  \item Perform gravity-resisted and gravity-eliminated testing for:
  \begin{itemize}
    \item \textbf{Hip Flexion:} Iliopsoas, Rectus Femoris, Sartorius
    \item \textbf{Hip Extension:} Gluteus Maximus, Hamstrings
  \end{itemize}
  \item Use proper stabilization and palpate for muscle activation.
  \item Record observed substitutions (e.g., lumbar lordosis with flexion; trunk extension with weak gluteals).
  \item Document innervation and nerve root(s) for each muscle tested.
\end{itemize}
\end{labactivitybox}

\begin{labactivitybox}
\textbf{Activity: MMT – Hip Abduction and Adduction}

\begin{itemize}
  \item Test:
  \begin{itemize}
    \item \textbf{Hip Abduction:} Gluteus Medius, Gluteus Minimus
    \item \textbf{Hip Adduction:} Adductor Longus, Adductor Magnus
  \end{itemize}
  \item Stabilize contralateral pelvis and avoid rolling or hip flexion substitutions.
  \item Compare side-to-side strength and assess symmetry.
  \item Match muscles tested with their peripheral nerves and spinal roots.
\end{itemize}
\end{labactivitybox}

\begin{labactivitybox}
\textbf{Activity: MMT – Hip Internal and External Rotation}

\begin{itemize}
  \item Test:
  \begin{itemize}
    \item \textbf{Internal Rotation:} Tensor Fascia Latae, Gluteus Medius (anterior)
    \item \textbf{External Rotation:} Piriformis, Obturator Externus
  \end{itemize}
  \item Observe for substitutions (e.g., pelvic hiking or leaning).
  \item Use visual landmarks to align the femur during rotation assessment.
  \item Identify innervation and root level for each muscle.
\end{itemize}
\end{labactivitybox}

\subsection*{MLT – Pelvic and Hip Muscles}

\begin{labactivitybox}
\textbf{Activity: MLT – Hamstrings (Straight Leg Raise)}

\begin{itemize}
  \item This repeats the SLR test introduced in Week 1, now used in context with pelvic stabilization and neuromechanical interpretation.
  \item Perform a passive straight leg raise (SLR) to assess hamstring extensibility.
  \item Maintain neutral hip rotation and stabilize contralateral limb.
  \item Record \textbf{R1} (first resistance) and \textbf{R2} (max available range) in degrees or descriptive terms.
  \item Note any compensations (e.g., posterior pelvic tilt, lumbar flexion).
  \item Classify length as Normal, Limited, or Excessive.
\end{itemize}
\end{labactivitybox}

\begin{labactivitybox}
\textbf{Activity: MLT – Hip Flexors (Modified Thomas Test)}

\begin{itemize}
  \item Use the modified Thomas test to assess length of:
  \begin{itemize}
    \item \textbf{Iliopsoas:} look for hip flexion (thigh not reaching table)
    \item \textbf{Rectus Femoris:} assess knee flexion angle
    \item \textbf{Tensor Fascia Latae (TFL):} assess hip abduction and external rotation
  \end{itemize}
  \item Document R1 and R2 with attention to pelvic tilt or lumbar extension.
  \item Interpret findings for each muscle based on limb position and movement.
\end{itemize}
\end{labactivitybox}

\begin{labactivitybox}
\textbf{Activity: MLT – TFL/IT Band (Ober’s Test)}

\begin{itemize}
  \item Position patient in sidelying with pelvis stabilized and bottom leg flexed.
  \item Passively abduct and extend the top leg, then allow it to lower toward the table.
  \item Assess for horizontal plane drop—restricted descent indicates tight TFL/IT band.
  \item Observe for pelvic hiking or trunk rotation.
  \item Record R1, R2, and classify length as Normal, Limited, or Excessive.
\end{itemize}
\end{labactivitybox}

\section*{Neurovascular RUSI}

The femoral artery is the primary vascular conduit to the lower extremity, emerging beneath the inguinal ligament and traveling through the femoral triangle alongside the femoral vein and femoral nerve. Given its superficial position in this region, it is easily visualized with ultrasound and serves as a useful site to practice probe stabilization, vessel identification, and tissue layer orientation.

Ultrasound imaging in this area allows visualization of all three components of the femoral neurovascular bundle:
\begin{itemize}
  \item \textbf{Femoral artery:} Pulsatile, anechoic, non-compressible.
  \item \textbf{Femoral vein:} Compressible, anechoic, medial to artery.
  \item \textbf{Femoral nerve:} Hyperechoic fascicular structure, lateral to artery.
\end{itemize}

\textbf{Note:} This activity introduces foundational RUSI skills for the anterior thigh. Additional RUSI applications—including quadriceps and joint-level imaging—will be covered in Week 6 during the knee-focused lab.

\begin{labactivitybox}
\textbf{Activity: RUSI – Femoral Neurovascular Bundle}

\begin{itemize}
  \item Place the probe in \textbf{SAX} orientation just distal to the inguinal ligament.
  \item Identify the femoral artery, vein, and nerve in cross-section.
  \item Use compression to differentiate vein from artery.
  \item Adjust depth and gain for optimal fascial contrast.
  \item Record the arrangement of the bundle: \textbf{NAV} (Nerve, Artery, Vein) from lateral to medial.
\end{itemize}
\end{labactivitybox}

\section*{Clinical Integration}

\begin{itemize}
  \item How did SFMA Top Tier patterns (e.g., MSF, MSE, MSR, SLS, ADS) help identify whether a patient’s issue was mobility or motor control driven?
  \item What neuromechanical relationships exist between the pelvis, hip, and lumbar spine? How might hip muscle weakness contribute to lumbopelvic dysfunction?
  \item Where did you observe altered length–tension relationships or compensations during MMT and MLT?
  \item How might asymmetries in hip control impact distal function (e.g., knee mechanics, gait)?
  \item How did RUSI of the femoral neurovascular bundle reinforce your understanding of regional anatomy and clinical safety considerations?
\end{itemize}

\section*{Next Steps}

\begin{itemize}
  \item Review the movement–muscle–nerve–root tables for all six primary hip motions.
  \item Practice MMT for hip flexion, extension, abduction, adduction, and rotation, with attention to stabilization and compensation.
  \item Repeat MLT for the hamstrings (SLR), rectus femoris (Modified Thomas), and TFL/IT band (Ober).
  \item Practice RUSI scanning of the femoral neurovascular bundle; reinforce vessel identification and probe orientation.
  \item Preview Week 6 content on the knee, with a focus on eccentric hamstring control, effusion imaging, and quadriceps testing.
\end{itemize}

\section*{Checklist Summary}

\begin{itemize}
  \item [$\square$] Interpreted SFMA Top Tier patterns (MSF, MSE, MSR, SLS, ADS) with attention to pelvic-hip coordination and regional interdependence.
  \item [$\square$] Performed gravity-resisted and gravity-eliminated MMT for major hip movements: flexion, extension, abduction, adduction, internal rotation, external rotation.
  \item [$\square$] Identified peripheral nerve and spinal root(s) for each primary muscle tested in MMT.
  \item [$\square$] Repeated the Straight Leg Raise (SLR) to assess hamstring extensibility and compare to Week 1 performance.
  \item [$\square$] Performed the Modified Thomas Test to assess iliopsoas and rectus femoris length using R1/R2 and compensatory pattern recognition.
  \item [$\square$] Conducted the Ober test for TFL/IT band extensibility with appropriate stabilization.
  \item [$\square$] Visualized the femoral neurovascular bundle using ultrasound and confirmed anatomical landmarks.
\end{itemize}
